% Verbatim historical narrative from the frozen LT v0.1-rc6 authority.
% Fragment, not a standalone paper. Mathematical proof text remains in reader.tex.
\subsection{Reconstruction boundary}
Revision 2 retains the first nonzero-angle prime family and
the final thirteen-sector tail route of checkpoint v38.
It expands the real-cutoff jet recurrence, separates the
radial-center and protected-disk arguments, and exposes the
hybrid measure and constant choices. These additions are
mathematical exposition to audit, not an external validation.

Revision 3 corrects the exterior convergence paragraph in
Appendix~\ref{app:pfaffian}: global inverse products are bounded
exponentially in \(N\), rather than independently of \(N\).
The displayed Pfaffian--factorial bound absorbs this cost, and
the compact-set Cauchy estimate makes the differentiated
normal convergence explicit. This revision leaves the main
claim, boundary-point family, and tail estimates unchanged.

Revision 4 corrects the one-body argument in
Lemma~\ref{lem:originaldisk}: the branch-center displacement
on a complex parameter disk includes \(|s-\se|\).
The two-region estimate now retains this perturbation and
proves the same occupancy-independent integrable majorant.
No contour, particle-number window, or main claim is changed.

Revision 5 incorporates the explicitly identified first-term gluing
supplement from the subsequent full-chain audit. Section~\ref{sec:first}
now constructs a \(y\)-only angular partition before taking any
\(x\) residues, replaces its smooth mean cutoffs only across separated
edge regions, and gives a quantitative finite-cover estimate for the
complement. These additions connect the local residue calculation to
the whole double torus. They are new exposition and supporting
derivations relative to revision 4, not a claim that revision 4 already
contained this construction. The tail estimates and the main target
are unchanged; the audit's other verification labels are not imported
as premises.

Revision 6 makes four further interfaces explicit following a
separate audit of revision 5: compact upper supports and branch
plateaus; a periodic left/right partition whose extra seam lies
where all relevant weighted jets vanish; Cauchy's inequality
applied to the complete holomorphic \(F_N\); and exponential-scale
inclusions imposed at every finite tail order. The related negative
anchor scale is chosen with the same uniform inclusion. These are
explicit choices and local justifications within the existing proof
route. The model, theorem, point family, derivative range, and
particle-number windows are unchanged. No audit verdict is used
as a mathematical premise.

Revision 7 corrects the logarithmic expansion of the microcore
majorant by displaying its negative \(N^2\log N/2\) term. It also
records the angle and power conventions required for a comparison
with the review amplitude, and explicitly includes the unwrapped
phase lengths and period counts in near-equality coarea estimates.
The physical-state convention is stated at the outset; the current
E1 integration specifies the \(+\) boundary-condition state. These are local corrections and
clarifications; the two main estimates, point family, derivative
range, and particle windows are unchanged. The received audit's
failure to find a counterexample is not an additional proof premise.

Revision 8 supplies the missing historical and methodological
attribution, including the diagonal noncancellation argument and the
upstream correlation representation. It distinguishes the proposed
bulk estimates from the existing proof architecture. The endpoint
wording in Lemma~\ref{lem:contraction} now asserts strictness only
within groups and explicitly permits cross-group equality. Inverse
temperature and the inner branch cutoff have distinct symbols from
the geometric angle \(\beta\). The selected family, two central
estimates, derivative range, and particle windows are unchanged.
The separate revision record states the checks actually performed;
earlier audit verdicts are attached to their earlier manuscript
versions and are not automatically new audits of revision 8.

Earlier routes involving a jointly absolute mean/shape
limit, a multiplier obtained by dividing by the original
Vandermonde, a central connector contour, or a degenerate
four-particle endpoint amplitude are not assumptions of
this manuscript. The separate source-coverage record names
their replacements in v38. A historical file being
present in the checkpoint is not by itself a reason to
include its superseded argument in the proof.

The accompanying release records identify file hashes, the active-source
mapping, and the edits from the preceding reconstruction.
Those checks establish provenance and textual coverage.
They do not establish truth of any proposition, settle
literature priority, or constitute a completed independent
adversarial audit.

% END SOURCE

\subsection*{Historical rc4 W2 repair}
Version 0.1-rc4 adds the explicit support geometry,
local root-continuation and complete-pair transfer arguments,
reserved disk slack, tau-dependent protected-disk restrictions,
and coupled-occupancy product-rule count. Lemma~\ref{lem:originaldisk}
now explicitly confines its pair estimates to \(K\) and named
Stokes supports. The theorem, point family, derivative range and
particle windows are unchanged. These are new scoped proof repairs,
not arguments attributed retrospectively to rc3. The supplied
audits and finite diagnostics were treated as untrusted evidence.
This candidate is prepared for public review. The subsequent E1 integration
checks the primary-source interface, specifies the \(+\) state, and
reconciles the published full-site coefficient with the existing
offsite formulas. It changes no W1/W2 estimate or theorem statement.
That analytic release did not claim human, peer, or proof-assistant certification.

\subsection*{Historical rc3 status and AI assistance}
Version 0.1-rc3 incorporates the W1 proof-interface repairs described in this revision record.
It adds the current support partition and deformed branch geometry,
recounts the remaining pair factors after differentiation, corrects
the near-collision slope cost, and expands the constant dependencies.
These substantive proof repairs received a scoped AI analytic audit;
the historical rc2 records retain their original inputs and verdicts.
Three problem-only reconstruction runs returned partial results;
none reconstructed the complete natural-boundary proof. One
method-informed run claimed the full tail estimate while explicitly
leaving its whole-integral first singular term (node I3) unresolved.
Its tail claim has not received a complete independent verification
by the receiving workflow. The different point families used in
some runs do not allow their conclusions to be combined silently.
The repository separates these run claims from manuscript claims,
version-specific audit reports, and finite computational checks.

Generative AI was used extensively in research exploration,
mathematical derivation, manuscript writing, code, internal audits,
and reconstruction attempts. Dehao Lin is the named human author
and public contact. No specific division of mathematical contributions
or line-by-line human verification is asserted here.
An AI audit verdict is not proof certification. No independent human
expert certification is claimed. Formal verification has the statement
and trust boundaries specified for the current candidate below.
The author declares no additional funding or competing interests
for this project. Project text is offered under CC BY 4.0 to the extent
the author holds rights in it; project code is offered under the MIT
license. Cited third-party works retain their own rights.
The accompanying provenance record states
the limits of the available historical access evidence.

\subsection*{Rc5 internal regularity and formal scope}
Version 0.1-rc5 was paired with Lean candidate 0.3.0-rc1.
The original rc4 source, PDF and label index remain separate frozen
artifacts. The internal lower-order proof removes E3 from the canonical
FIRST and natural-boundary theorem statements and their proof dependency
chain. The constructed normalized exterior series has no external physics
input; its identification with the physical susceptibility uses E1 and E2.
The stronger selected-contour pair bound was already proved by the earlier
Lean source; this revision makes that strength explicit in the manuscript.
The closed-form amplitude evaluation and finite-prime covering calculation
are analytic manuscript results, not new Lean exports of those formulas.
Lean proves the actual selected-point real \(C^k\) obstruction with the tail
internally discharged. The all-arc quantitative covering is not formalized.

\subsection*{Rc6 complete-amplitude comparison and formal scope}
The current manuscript is 0.1-rc6, paired with Lean candidate 0.3.0-rc2.
Proposition~\ref{prop:amplitudecomparison} identifies the full-site
observable from the review's cited integral, fixes the exterior branch,
and proves the four-chart conversion of its printed scalar. It also
distinguishes the lower double-contour coefficient from the full offsite
lattice coefficient and includes both physical prefactors. The new
conversion and radial Taylor reconstruction are manuscript results;
the factorial and literature-amplitude identities are not new Lean exports.
The production Lean theorem sources and their mathematical dependency
chain are unchanged from the verified rc5/0.3.0-rc1 candidate. Updated
source correspondence and verification identities belong to this revision;
the preserved commit records the earlier candidate without reassigning its
receipts to new bytes.

Exact build, transitive-dependency, axiom, kernel-replay, and PDF receipts
are supplied separately with their candidate hashes. They do not extend
earlier AI audit verdicts to these new bytes. Dependency binary/source
correspondence remains trusted, and build and replay use the same Lean
kernel implementation. Neither these internal checks nor compilation is
independent human review or a literature-priority certificate.

